% GENERATED FILE. Edit canonical lessons, then run npm run generate:books.
% canonical-source-hash: fnv1a64:f964923ecf3493b3
% canonical-lessons: MR-C540-vima, MR-C540-jokhim, MR-C540-hami, MR-C540-hapta, MR-C540-parvadne

\chapter{Insurance and Risk}
\label{ch:mr-insurance-and-risk}
\hlchaptermodality{540}

\begin{chapteropening}
\textbf{By the end of this chapter:} \emph{I can talk about insurance, risks and what I can afford.}

Complete the last lesson of chapter 540: I can talk about insurance, risks and what I can afford.
\end{chapteropening}

\section[\texorpdfstring{\mr{विमा} (vimā)}{विमा (vimā)}]{\mr{विमा} (vimā) — insurance}

\label{lesson:MR-C540-vima}



\begin{quote}
Before the new one: say the Marathi for a film director, then the Marathi for to applaud.
\end{quote}

\subsection*{You'll want to know: \mr{विमा}}
\textbf{\mr{विमा}} — \emph{vimā} — \textquotedblleft{}insurance\textquotedblright{}.

\begin{quote}
\mr{जर} \mr{घराचा} \mr{विमा} \mr{असेल} \mr{तर} \mr{नुकसान} \mr{भरून} \mr{मिळतं}.
\end{quote}

\begin{quote}
\emph{jar gharācā vimā asel tar nuksān bharūn miḷtaṁ.}
\end{quote}

\begin{quote}
If the house is insured, the damage is made good.
\end{quote}

Masculine.

\mr{विमा} \mr{उतरवणे}, to take out insurance.

\begin{cousinweb}
\textbf{\mr{विमा}} (\emph{vimā}) is a loan from Persian \emph{bīma}, \textquotedblleft{}insurance, fear\textquotedblright{}.
\end{cousinweb}

\begin{grammarlens}[title={What you've built}]
One more word for saying what happens if.
\end{grammarlens}

\subsection*{Your turn}
\begin{itemize}
\raggedright
  \item \emph{Say it:} \emph{vimā}
  \item \emph{Say it:} \emph{vimā}, once more
  \item \emph{Say it:} \emph{jar gharācā vimā asel tar nuksān bharūn miḷtaṁ.} — If the house is insured, the damage is made good.
  \item \emph{Recall:} say the Marathi for a film director, then the Marathi for to applaud, then say \emph{vimā} again
\end{itemize}

\begin{tcolorbox}[breakable,colback=teal!4,colframe=teal!35!black,title={Before you move on}]
What does \mr{विमा} mean? (\textquotedblleft{}Insurance\textquotedblright{}.) Say it once more.
\end{tcolorbox}
\section[\texorpdfstring{\mr{जोखीम} (jokhīm)}{जोखीम (jokhīm)}]{\mr{जोखीम} (jokhīm) — a risk}

\label{lesson:MR-C540-jokhim}



\begin{quote}
Before the new one: say the Marathi for to applaud, then the Marathi for insurance.
\end{quote}

\subsection*{You'll want to know: \mr{जोखीम}}
\textbf{\mr{जोखीम}} — \emph{jokhīm} — \textquotedblleft{}a risk\textquotedblright{}.

\begin{quote}
\mr{कोणताही} \mr{व्यवसाय} \mr{सुरू} \mr{करताना} \mr{थोडी} \mr{जोखीम} \mr{असतेच}.
\end{quote}

\begin{quote}
\emph{koṇtāhī vyavsāy surū kartānā thoḍī jokhīm astec.}
\end{quote}

\begin{quote}
Starting any business always carries some risk.
\end{quote}

Feminine.

\mr{जोखीम} \mr{घेणे}, to take a risk.

\begin{cousinweb}
\textbf{\mr{जोखीम}} (\emph{jokhīm}) is a tadbhava: CDIAL lists Marathi \emph{jokhīm}, \textquotedblleft{}risk, venture\textquotedblright{}, under Sanskrit \emph{yōgakṣēma}, \textquotedblleft{}the safe keeping of one's property\textquotedblright{} (CDIAL 10527).
\end{cousinweb}

\begin{grammarlens}[title={What you've built}]
One more word for saying what happens if.
\end{grammarlens}

\subsection*{Your turn}
\begin{itemize}
\raggedright
  \item \emph{Say it:} \emph{jokhīm}
  \item \emph{Say it:} \emph{jokhīm}, once more
  \item \emph{Say it:} \emph{koṇtāhī vyavsāy surū kartānā thoḍī jokhīm astec.} — Starting any business always carries some risk.
  \item \emph{Recall:} say the Marathi for to applaud, then the Marathi for insurance, then say \emph{jokhīm} again
\end{itemize}

\begin{tcolorbox}[breakable,colback=teal!4,colframe=teal!35!black,title={Before you move on}]
What does \mr{जोखीम} mean? (\textquotedblleft{}A risk\textquotedblright{}.) Say it once more.
\end{tcolorbox}
\section[\texorpdfstring{\mr{हमी} (hamī)}{हमी (hamī)}]{\mr{हमी} (hamī) — a guarantee}

\label{lesson:MR-C540-hami}



\begin{quote}
Before the new one: say the Marathi for insurance, then the Marathi for a risk.
\end{quote}

\subsection*{You'll want to know: \mr{हमी}}
\textbf{\mr{हमी}} — \emph{hamī} — \textquotedblleft{}a guarantee\textquotedblright{}.

\begin{quote}
\mr{या} \mr{पंख्यावर} \mr{दोन} \mr{वर्षांची} \mr{हमी} \mr{आहे}.
\end{quote}

\begin{quote}
\emph{yā paṅkhyāvar don varṣāñcī hamī āhe.}
\end{quote}

\begin{quote}
This fan has a two-year guarantee.
\end{quote}

Feminine.

\mr{हमी} \mr{देणे}, to guarantee.

\begin{cousinweb}
\textbf{\mr{हमी}} (\emph{hamī}) is the word for a guarantee or pledge; its origin is not settled.
\end{cousinweb}

\begin{grammarlens}[title={What you've built}]
One more word for saying what happens if.
\end{grammarlens}

\subsection*{Your turn}
\begin{itemize}
\raggedright
  \item \emph{Say it:} \emph{hamī}
  \item \emph{Say it:} \emph{hamī}, once more
  \item \emph{Say it:} \emph{yā paṅkhyāvar don varṣāñcī hamī āhe.} — This fan has a two-year guarantee.
  \item \emph{Recall:} say the Marathi for insurance, then the Marathi for a risk, then say \emph{hamī} again
\end{itemize}

\begin{tcolorbox}[breakable,colback=teal!4,colframe=teal!35!black,title={Before you move on}]
What does \mr{हमी} mean? (\textquotedblleft{}A guarantee\textquotedblright{}.) Say it once more.
\end{tcolorbox}
\section[\texorpdfstring{\mr{हप्ता} (haptā)}{हप्ता (haptā)}]{\mr{हप्ता} (haptā) — an instalment; a premium}

\label{lesson:MR-C540-hapta}



\begin{quote}
Before the new one: say the Marathi for a risk, then the Marathi for a guarantee.
\end{quote}

\subsection*{You'll want to know: \mr{हप्ता}}
\textbf{\mr{हप्ता}} — \emph{haptā} — \textquotedblleft{}an instalment; a premium\textquotedblright{}.

\begin{quote}
\mr{विम्याचा} \mr{हप्ता} \mr{दर} \mr{महिन्याला} \mr{भरावा} \mr{लागतो}.
\end{quote}

\begin{quote}
\emph{vimyācā haptā dar mahinyālā bharāvā lāgto.}
\end{quote}

\begin{quote}
The insurance premium has to be paid every month.
\end{quote}

Masculine.

\mr{हप्ता} \mr{भरणे}, to pay an instalment.

\begin{cousinweb}
\textbf{\mr{हप्ता}} (\emph{haptā}) is a loan from Persian \emph{hafta}, \textquotedblleft{}week\textquotedblright{}: a payment made at a set interval.
\end{cousinweb}

\begin{grammarlens}[title={What you've built}]
One more word for saying what happens if.
\end{grammarlens}

\subsection*{Your turn}
\begin{itemize}
\raggedright
  \item \emph{Say it:} \emph{haptā}
  \item \emph{Say it:} \emph{haptā}, once more
  \item \emph{Say it:} \emph{vimyācā haptā dar mahinyālā bharāvā lāgto.} — The insurance premium has to be paid every month.
  \item \emph{Recall:} say the Marathi for a risk, then the Marathi for a guarantee, then say \emph{haptā} again
\end{itemize}

\begin{tcolorbox}[breakable,colback=teal!4,colframe=teal!35!black,title={Before you move on}]
What does \mr{हप्ता} mean? (\textquotedblleft{}An instalment; a premium\textquotedblright{}.) Say it once more.
\end{tcolorbox}
\section[\texorpdfstring{\mr{परवडणे} (parvaḍṇe)}{परवडणे (parvaḍṇe)}]{\mr{परवडणे} (parvaḍṇe) — to be affordable}

\label{lesson:MR-C540-parvadne}



\begin{quote}
Before the new one: say the Marathi for a guarantee, then the Marathi for an instalment.
\end{quote}

\subsection*{You'll want to know: \mr{परवडणे}}
\textbf{\mr{परवडणे}} — \emph{parvaḍṇe} — \textquotedblleft{}to be affordable\textquotedblright{}.

\begin{quote}
\mr{इतकी} \mr{महाग} \mr{गाडी} \mr{आम्हाला} \mr{परवडणार} \mr{नाही}.
\end{quote}

\begin{quote}
\emph{itkī mahāg gāḍī āmhālā parvaḍṇār nāhī.}
\end{quote}

\begin{quote}
We cannot afford such an expensive car.
\end{quote}

Past (to him …): \mr{त्याला} \mr{परवडलं}, \emph{tyālā parvaḍlaṁ}.

The person takes -\mr{ला}: \mr{मला} \mr{परवडतं}, I can afford it.

\begin{cousinweb}
\textbf{\mr{परवडणे}} (\emph{parvaḍṇe}) is a Marathi verb with the one who pays in the dative; its origin is not settled.
\end{cousinweb}

\begin{grammarlens}[title={What you've built}]
One more word for saying what happens if.
\end{grammarlens}

\subsection*{Your turn}
\begin{itemize}
\raggedright
  \item \emph{Say it:} \emph{parvaḍṇe}
  \item \emph{Say it:} \emph{parvaḍṇe}, once more
  \item \emph{Say it:} \emph{itkī mahāg gāḍī āmhālā parvaḍṇār nāhī.} — We cannot afford such an expensive car.
  \item \emph{Recall:} say the Marathi for a guarantee, then the Marathi for an instalment, then say \emph{parvaḍṇe} again
\end{itemize}

\begin{tcolorbox}[breakable,colback=teal!4,colframe=teal!35!black,title={Before you move on}]
What does \mr{परवडणे} mean? (\textquotedblleft{}To be affordable\textquotedblright{}.) Say it once more.
\end{tcolorbox}
